\section{讨论与效度威胁}
\sectionclaim{本系统的效度取决于证据构念、模型核查、查询分布和小样本推断四个因素。}

\textbf{构念效度由冻结证据组与问题语义共同决定。}备选摘录经原文及合格块双重验证，但来源审核仍可能遗漏等价表达。IOC 清单的部分字段来自面板描述，“启动上报”与“界面展示”的语义边界值得独立复核；本文完整分数始终沿用冻结 v2 构念，文献出版年与本地版本标签分别保留。

\textbf{内部效度受到核查器与生成器共享模型的影响。}33 个模型提取要求与 35 个冻结证据组采用不同粒度，单个宽泛要求可以接受部分字段清单。IOC 实例证明存在过度完整标记，双人盲审将用于估计其发生率；数字与引用 ID 检查只验证可追溯性，语义支持仍由核查过程判定。

\textbf{外部效度受开发查询与静态语料分布约束。}21 条可回答问题已用于方案比较，3 条跨语言目标为重叠小切片；领域负例主要覆盖时间性缺口。独立保留集、复杂 PDF 结构和运行数据接口是后续扩展方向；\texttt{ds-1017}仅完成代表性验证，本文 32 问服务指标固定归属\texttt{qwen3.8-Flash-Next}。

\textbf{结论效度通过配对效应与原始轨迹共同界定。}两条不一致样本给出\(p=0.500\)，差值区间包含零，当前主张是候选互补机制与固定问题集上的完整覆盖改善。bootstrap 区间描述经验查询重采样波动，方法筛选、相关查询和多重比较影响总体推广；计时采用单次稳态记录，多查询规划耗时及同环境单查询对照缺失，跨 CPU/GPU 或模型版本的因果加速比留待受控测量。
